\begin{abstract}
Scientific-agent benchmarks rarely test whether an agent can allocate a fixed set of
physical units, use evidence that arrives during long-running trials, and commit to a
final operating recommendation whose experimental cost is paid in the objective itself.
We introduce \slowlab, an executable controlled-environment horticulture benchmark with
split-plot controls, capped parallel capacity, correlated noise, irreversible 210-day
crop cycles, time-appropriate interim observations, and separate design and
recommendation actions. We evaluate final and cumulative regret and use a common-history
Bayes-risk decomposition to separate what an executed design supports from how a reader
uses that evidence. Across 928 episodes under a frozen protocol, neither design nor
inference aids significantly changed final regret for Luna on 56 paired sites, and an
inference aid did not improve Sol on 24 further sites. A fresh 384-episode extension found
the same null result at twice the observation noise and a harmful design-aid effect at half
noise ($\Delta=0.0113$, 95\% CI $[0.0053,0.0184]$, adjusted $p=0.0096$). In a prospective
post-primary 288-episode same-site extension, the inference aid showed no detectable change for Luna, was
directionally favorable but uncertain for DeepSeek V4.1 Flash, and reduced Qwen 3.8 27B
final regret by $0.0146$ (95\% CI $[0.0091,0.0201]$; BH $p<0.001$), showing that tool value
is model dependent under the tested interfaces.
\end{abstract}
